% IMPORTANT: This LaTeX file is exported by ATHENA. Please do NOT change ANY ATHENA-DATA record to avoid loss of ATHENA interoperability.
% ATHENA is free software. To learn more about ATHENA please visit https://athena.evalisk.org/
% ATHENA-DATA cmd="version" val=("0.9")
% ATHENA-DATA cmd="aux" val=("document_style", "20", "PHR1cGxlfGdlbmVyaWM+")
% ATHENA-DATA cmd="aux" val=("document_initial", "208", "PFxjb2xsZWN0aW9uPgogIDxhc3NvY2lhdGV8Zm9udHx0eXBld3JpdGVyPUpldEJyYWlucyBNb25vLFRlWCBHeXJlIFBhZ2VsbGE+CiAgPGFzc29jaWF0ZXxmb250LWZhbWlseXxybT4KICA8YXNzb2NpYXRlfHBhZ2UtbWVkaXVtfGF1dG9tYXRpYz4KPC9jb2xsZWN0aW9uPg==")
\documentclass{article}
\long\def\INLINE_COMMENT#1{}
\usepackage[english]{babel}
\usepackage[utf8]{inputenc}
\usepackage{amsmath,graphicx,hyperref,alltt,theorem}

%%%%%%%%%% Start TeXmacs macros
\newcommand{\ATHENA}{\textbf{\kern.07emA\kern-.07em\raisebox{.25ex}{\scalebox{.7}{T}}\kern-.05emH\kern.05em\raisebox{-.25ex}{\scalebox{.7}{E}}NA}}
\newcommand{\tmverbatim}[1]{\text{{\ttfamily{#1}}}}
\newenvironment{tmcode}[1][]{\begin{alltt} }{\end{alltt}}
{\theorembodyfont{\rmfamily}\newtheorem{remark}{Remark}}
%%%%%%%%%% End TeXmacs macros

\begin{document}

\title{{\ATHENA} - Compiling}

\date{}

\maketitle

\INLINE_COMMENT{ATHENA-DATA cmd="object" val=("32", "PGxhYmVsfEgxIEludHJvZHVjdGlvbj4=")}

\section{Introduction}

This document describes how to build {\ATHENA} from source.

The supported build target is x86\_64 GNU/Linux with Qt 6 and Wayland. The
reference development environment is openSUSE Tumbleweed snapshot 20260812,
KDE Plasma/KWin 6.7.4 on Wayland, and Qt 6.11.1. The Tumbleweed build with
Intel oneAPI \tmverbatim{icx}/\tmverbatim{icpx} is the configuration currently
expected to compile without source changes.

Intel oneAPI C++ (\tmverbatim{icpx}) is the recommended compiler. Install the
current Intel oneAPI DPC++/C++ Compiler using the instructions on the
\href{https://www.intel.com/content/www/us/en/developer/tools/oneapi/dpc-compiler-download.html}{Intel
compiler download and installation page}.

\INLINE_COMMENT{ATHENA-DATA cmd="object" val=("48", "PGxhYmVsfEgxIEJ1aWxkaW5nIG9uIHg2NCBHTlUvTGludXg+")}

\section{Building {\ATHENA} on x64 GNU/Linux}

To build {\ATHENA} for direct use, follow sections 2.1 through 2.4. To build
an RPM, DEB, or AppImage, obtain the source as in 2.1 and then use 2.5.

\INLINE_COMMENT{ATHENA-DATA cmd="object" val=("44", "PGxhYmVsfEgyIE9idGFpbmluZyB0aGUgc291cmNlPg==")}

\subsection{Obtaining the source}

\begin{tmcode}
git clone https://github.com/NuaptanEvalisk/ATHENA.git
cd ATHENA
\end{tmcode}
Checkout the branch or tag that you intend to build.

\INLINE_COMMENT{ATHENA-DATA cmd="object" val=("48", "PGxhYmVsfEgyIEluc3RhbGxpbmcgZGVwZW5kZW5jaWVzPg==")}

\subsection{Installing dependencies}

Install the corresponding development packages for your distribution. On
current openSUSE Tumbleweed, the following is an example:
\begin{tmcode}
sudo zypper install \textbackslash
  git cmake ninja pkgconf-pkg-config patchelf \textbackslash
  autoconf automake libtool make makeinfo gperf bison flex \textbackslash
  qt6-base-devel qt6-svg-devel qt6-qt5compat-devel \textbackslash
  qt6-wayland-devel qt6-waylandclient-devel \textbackslash
  kf6-kio-devel kf6-kconfigwidgets-devel \textbackslash
  kf6-kcompletion-devel kf6-syntax-highlighting-devel \textbackslash
  boost-devel mimalloc-devel sqlite3-devel spdlog-devel hunspell-devel \textbackslash
  freetype2-devel fontconfig-devel harfbuzz-devel \textbackslash
  libpng16-devel libjpeg8-devel zlib-devel libzstd-devel \textbackslash
  libsodium-devel ImageMagick-devel libgnutls-devel gmp-devel \textbackslash
  libffi-devel libunistring-devel readline-devel ncurses-devel \textbackslash
  libcurl-devel ghostscript resvg-devel \textbackslash
  python3 python313-uv cargo rust
\end{tmcode}
\INLINE_COMMENT{ATHENA-DATA cmd="object" val=("156", "PGxhYmVsfHJlbWFyazpBIHN5c3RlbSBHdWlsZSBpbnN0YWxsYXRpb24gaXMgbmVpdGhlciBuZWVkZWQgbm9yIHVzZWQgQ01ha2UKYnVpbGRzIHRoZSB2ZW5kb3JlZCBHdWlsZSAzIHJ1bnRpbWUgYXV0IHs+")}

\begin{remark}
  A system Guile installation is neither needed nor used. CMake builds the
  vendored {\ATHENA} Guile 3 runtime automatically. Do not configure the build
  against system Guile.
\end{remark}

\INLINE_COMMENT{ATHENA-DATA cmd="object" val=("156", "PGxhYmVsfHJlbWFyazpBIHN5c3RlbSBHdWlsZSBpbnN0YWxsYXRpb24gaXMgbmVpdGhlciBuZWVkZWQgbm9yIHVzZWQgQ01ha2UKYnVpbGRzIHRoZSB2ZW5kb3JlZCBHdWlsZSAzIHJ1bnRpbWUgYXV0IH0+")}

\INLINE_COMMENT{ATHENA-DATA cmd="object" val=("36", "PGxhYmVsfEgzIENvbXBpbGluZyBsaWJ0Y2M+")}

\subsubsection*{Compiling libtcc}

If no usable libtcc development package is available, build TinyCC from
source:
\begin{tmcode}
git clone https://repo.or.cz/tinycc.git
cd tinycc
CC=icx ./configure
make -j16
\end{tmcode}
Pass the resulting header and library to the {\ATHENA} CMake configuration:
\begin{tmcode}
-DTCC_INCLUDE_DIR=/path/to/tinycc \textbackslash
-DTCC_LIBRARY=/path/to/tinycc/libtcc.a
\end{tmcode}
\INLINE_COMMENT{ATHENA-DATA cmd="object" val=("40", "PGxhYmVsfEgzIENvbXBpbGluZyBsbGFtYS5jcHA+")}

\subsubsection*{Compiling llama.cpp}

{\ATHENA} currently builds against the llama.cpp revision used by its
container toolchain: \tmverbatim{66c4f9ded01b29d9120255be1ed8d5835bcbb51d}. A
CPU build can be produced with:
\begin{tmcode}
git clone https://github.com/ggml-org/llama.cpp.git
cd llama.cpp
git checkout 66c4f9ded01b29d9120255be1ed8d5835bcbb51d
source /opt/intel/oneapi/setvars.sh
cmake -S . -B build -G Ninja \textbackslash
  -DCMAKE_BUILD_TYPE=Release \textbackslash
  -DCMAKE_C_COMPILER=icx \textbackslash
  -DCMAKE_CXX_COMPILER=icpx \textbackslash
  -DBUILD_SHARED_LIBS=ON \textbackslash
  -DLLAMA_BUILD_COMMON=ON \textbackslash
  -DLLAMA_BUILD_EXAMPLES=OFF \textbackslash
  -DLLAMA_BUILD_TESTS=OFF \textbackslash
  -DLLAMA_BUILD_TOOLS=OFF
cmake --build build -j16
\end{tmcode}
For an Intel SYCL-enabled build, configure llama.cpp with the additional
backend options:
\begin{tmcode}
cmake -S . -B build -G Ninja \textbackslash
  -DCMAKE_BUILD_TYPE=Release \textbackslash
  -DCMAKE_C_COMPILER=icx \textbackslash
  -DCMAKE_CXX_COMPILER=icpx \textbackslash
  -DBUILD_SHARED_LIBS=ON \textbackslash
  -DGGML_SYCL=ON \textbackslash
  -DGGML_SYCL_TARGET=INTEL \textbackslash
  -DGGML_SYCL_DNN=ON \textbackslash
  -DGGML_SYCL_GRAPH=ON \textbackslash
  -DLLAMA_BUILD_COMMON=ON \textbackslash
  -DLLAMA_BUILD_EXAMPLES=OFF \textbackslash
  -DLLAMA_BUILD_TESTS=OFF \textbackslash
  -DLLAMA_BUILD_TOOLS=OFF
cmake --build build -j16
\end{tmcode}
\INLINE_COMMENT{ATHENA-DATA cmd="object" val=("36", "PGxhYmVsfEgyIENvbXBpbGluZyBBVEhFTkE+")}

\subsection{Compiling {\ATHENA}}

Load the oneAPI environment and choose the build parallelism. The
\tmverbatim{ATHENA\_BUILD\_JOBS} variable controls the private runtime builds;
\tmverbatim{-j} controls the outer CMake build.
\begin{tmcode}
source /opt/intel/oneapi/setvars.sh
JOBS=16
export ATHENA_BUILD_JOBS="$JOBS"
\end{tmcode}
For a Debug build in \tmverbatim{build\_qt6}:
\begin{tmcode}
CC=icx CXX=icpx cmake -S . -B build_qt6 -G Ninja \textbackslash
  -DATHENA_GUI=Qt6 \textbackslash
  -DCMAKE_BUILD_TYPE=Debug \textbackslash
  -DLLAMA_CPP_SOURCE_DIR=/path/to/llama.cpp
cmake --build build_qt6 -j"$JOBS"
\end{tmcode}
For a \tmverbatim{RelWithDebInfo} build in \tmverbatim{build\_rel}:
\begin{tmcode}
CC=icx CXX=icpx cmake -S . -B build_rel -G Ninja \textbackslash
  -DATHENA_GUI=Qt6 \textbackslash
  -DCMAKE_BUILD_TYPE=RelWithDebInfo \textbackslash
  -DLLAMA_CPP_SOURCE_DIR=/path/to/llama.cpp
cmake --build build_rel -j"$JOBS"
\end{tmcode}
If libtcc was built in a non-system location, add the
\tmverbatim{TCC\_INCLUDE\_DIR} and \tmverbatim{TCC\_LIBRARY} arguments from
2.2 to the CMake configuration.

\INLINE_COMMENT{ATHENA-DATA cmd="object" val=("156", "PGxhYmVsfHJlbWFyazpUaGUgZmlyc3QgYnVpbGQgbWF5IHNwZW5kIGEgc3Vic3RhbnRpYWwgYW1vdW50IG9mIHRpbWUKYm9vdHN0cmFwcGluZyB0aGUgdmVuZG9yZWQgR3VpbGUgcnVudGltZSBUaGluIHs+")}

\begin{remark}
  The first build may spend a substantial amount of time bootstrapping the
  vendored Guile runtime. ThinLTO links may also take a long time.
\end{remark}

\INLINE_COMMENT{ATHENA-DATA cmd="object" val=("156", "PGxhYmVsfHJlbWFyazpUaGUgZmlyc3QgYnVpbGQgbWF5IHNwZW5kIGEgc3Vic3RhbnRpYWwgYW1vdW50IG9mIHRpbWUKYm9vdHN0cmFwcGluZyB0aGUgdmVuZG9yZWQgR3VpbGUgcnVudGltZSBUaGluIH0+")}

\INLINE_COMMENT{ATHENA-DATA cmd="object" val=("76", "PGxhYmVsfEgzIENvbXBpbGluZyBBVEhFTkEgd2l0aCBUaHJlYWRTYW5pdGl6ZXIgZW5hYmxlZD4=")}

\subsubsection*{Compiling {\ATHENA} with ThreadSanitizer enabled}

\begin{tmcode}
CC=icx CXX=icpx cmake -S . -B build_tsan -G Ninja \textbackslash
  -DATHENA_GUI=Qt6 \textbackslash
  -DCMAKE_BUILD_TYPE=Debug \textbackslash
  -DATHENA_ENABLE_TSAN=ON \textbackslash
  -DLLAMA_CPP_SOURCE_DIR=/path/to/llama.cpp
cmake --build build_tsan -j"$JOBS"
\end{tmcode}
\INLINE_COMMENT{ATHENA-DATA cmd="object" val=("156", "PGxhYmVsfHJlbWFyazpUbyBidWlsZCB0aGUgdGVzdCBzdWl0ZSBjb25maWd1cmUgd2l0aCBEQlVJTERURVNUU09OIFJ1biBpdAp3aXRoIGN0ZXN0IHRlc3RkaXIgYnVpbGRxdDYgb3V0cHV0b25mYWlsIHs+")}

\begin{remark}
  To build the test suite, configure with \tmverbatim{-DBUILD\_TESTS=ON}. Run
  it with \tmverbatim{ctest --test-dir build\_qt6 --output-on-failure} or the
  corresponding build directory.
\end{remark}

\INLINE_COMMENT{ATHENA-DATA cmd="object" val=("156", "PGxhYmVsfHJlbWFyazpUbyBidWlsZCB0aGUgdGVzdCBzdWl0ZSBjb25maWd1cmUgd2l0aCBEQlVJTERURVNUU09OIFJ1biBpdAp3aXRoIGN0ZXN0IHRlc3RkaXIgYnVpbGRxdDYgb3V0cHV0b25mYWlsIH0+")}

\INLINE_COMMENT{ATHENA-DATA cmd="object" val=("52", "PGxhYmVsfEgyIEFzc2VtYmxpbmcgdGhlIEFUSEVOQSB0cmVlPg==")}

\subsection{Assembling the {\ATHENA} tree}

The default build stages the private Guile runtime and compiled Scheme modules
into \tmverbatim{ATHENA/}. Assemble or refresh the remaining runtime files
with the outputs from the selected build directory:
\begin{tmcode}
BUILD=build_qt6
install -m755 "$BUILD/src/ATHENA.bin" ATHENA/bin/
install -m755 "$BUILD/src/ATHENA-Watchdog" ATHENA/bin/
install -m755 "$BUILD/src/athena-codex-bridge" ATHENA/bin/
install -m755 "$BUILD/materials-engine-cargo/release/athena-materials-engine" ATHENA/bin/
install -m755 "$BUILD/tools/athena-transmitter/athena-transmitter" ATHENA/bin/
install -m755 "$BUILD/tools/athena-web-server/athena-web-server" ATHENA/bin/
mkdir -p ATHENA/lib ATHENA/share/ATHENA/web
cp -a "$BUILD"/x64/lib/libqt6advanceddocking*.so* ATHENA/lib/
cp -a /path/to/llama.cpp/build/bin/libllama.so* ATHENA/lib/
cp -a /path/to/llama.cpp/build/bin/libggml*.so* ATHENA/lib/
cp -a tools/athena-web-server/web/. ATHENA/share/ATHENA/web/
\end{tmcode}
Use \tmverbatim{BUILD=build\_rel} for a RelWithDebInfo tree. For a SYCL
llama.cpp build, include the matching oneAPI runtime libraries when they are
not provided by the target system.

\INLINE_COMMENT{ATHENA-DATA cmd="object" val=("68", "PGxhYmVsfEgyIEJ1aWxkaW5nIFJQTSwgREVCLCBhbmQgQXBwSW1hZ2UgcGFja2FnZXM+")}

\subsection{Building RPM, DEB, and AppImage packages}

The container builder requires Podman and expects Intel oneAPI under
\tmverbatim{/opt/intel/oneapi}. To build the release flavor:
\begin{tmcode}
ATHENA_BUILD_JOBS=16 ATHENA_BUILD_FLAVORS=rel \textbackslash
  tools/container-build/build-athena-container.sh
\end{tmcode}
The principal outputs are:
\begin{itemize}
  \item \tmverbatim{container\_build/ATHENA-rel.AppImage}.

  \item
  \tmverbatim{container\_build/packages/ATHENA-VERSION-rel-linux-x86\_64.deb}.

  \item
  \tmverbatim{container\_build/packages/ATHENA-VERSION-rel-linux-x86\_64.rpm}.
\end{itemize}
Set \tmverbatim{ATHENA\_BUILD\_FLAVORS} to \tmverbatim{dev}, \tmverbatim{rel},
or both as required.

\INLINE_COMMENT{ATHENA-DATA cmd="object" val=("60", "PGxhYmVsfEgxIEJ1aWxkaW5nIEFUSEVOQSBvbiBvdGhlciBwbGF0Zm9ybXM+")}

\section{Building {\ATHENA} on other platforms}

\INLINE_COMMENT{ATHENA-DATA cmd="object" val=("52", "PGxhYmVsfEgyIEJ1aWxkaW5nIEFUSEVOQSBmb3IgV2luZG93cz4=")}

\subsection{Building {\ATHENA} for Windows}

Windows is not a tested or actively maintained target. A native Windows build
is theoretically possible with the necessary platform work.

The repository contains a MinGW-w64 cross-build pipeline under
\tmverbatim{tools/windows-deps}. It is currently incomplete because the
vendored, modified Guile 3.0.10 runtime has not been ported to the Windows
target. The pipeline must be adapted to build and package that private Guile
runtime and its BDW-GC dependency before it can produce a current {\ATHENA}
Windows build. The old Guile 1.8 path is obsolete.

\INLINE_COMMENT{ATHENA-DATA cmd="object" val=("92", "PGxhYmVsfEgyIEJ1aWxkaW5nIEFUSEVOQSBmb3IgbWFjT1MsIGFybTY0IExpbnV4LCBhbmQgb3RoZXIgdGFyZ2V0cz4=")}

\subsection{Building {\ATHENA} for macOS, arm64 Linux, and other targets}

macOS, arm64 GNU/Linux, and other platforms are theoretically possible with
appropriate source, dependency, and toolchain modifications. They are not
currently tested or supported build targets.

\INLINE_COMMENT{ATHENA-DATA cmd="object" val=("52", "PGxhYmVsfEgxIEdldHRpbmcgU3RhcnRlZCB3aXRoIEFUSEVOQT4=")}

\section{Getting Started with {\ATHENA}}

\INLINE_COMMENT{ATHENA-DATA cmd="object" val=("36", "PGxhYmVsfEgyIExhdW5jaGluZyBBVEhFTkE+")}

\subsection{Launching {\ATHENA}}

If you followed sections 2.1 through 2.4, start {\ATHENA} with:
\begin{tmcode}
cd ATHENA
./StartATHENA.qt6.wayland.sh
\end{tmcode}
To run through the watchdog:
\begin{tmcode}
./StartATHENA.qt6.wayland.sh --enable-watchdog
\end{tmcode}
If you followed 2.1 and 2.5, install the RPM or DEB with the package manager
of the target distribution, or run the AppImage directly.

\INLINE_COMMENT{ATHENA-DATA cmd="object" val=("156", "PGxhYmVsfHJlbWFyazpYMTFYQ0Igc3VwcG9ydCBpcyBub3QgYWN0aXZlbHkgbWFpbnRhaW5lZCBhbmQgbWF5IGNvbnRhaW4KZGVmZWN0cyBOYXRpdmUgV2F5bGFuZCBpcyB0aGUgbWFpbnRhaW5lZCBkIHs+")}

\begin{remark}
  X11/XCB support is not actively maintained and may contain defects. Native
  Wayland is the maintained desktop path.
\end{remark}

\INLINE_COMMENT{ATHENA-DATA cmd="object" val=("156", "PGxhYmVsfHJlbWFyazpYMTFYQ0Igc3VwcG9ydCBpcyBub3QgYWN0aXZlbHkgbWFpbnRhaW5lZCBhbmQgbWF5IGNvbnRhaW4KZGVmZWN0cyBOYXRpdmUgV2F5bGFuZCBpcyB0aGUgbWFpbnRhaW5lZCBkIH0+")}

\INLINE_COMMENT{ATHENA-DATA cmd="object" val=("156", "PGxhYmVsfHJlbWFyazpkb2NraW5nIHVzZXMgeGRndG9wbGV2ZWxkcmFndjEgS0RFIFBsYXNtYSA2MSBhbmQgbGF0ZXIgc3VwcG9ydAp0aGUgbWVyZ2VkIHByb3RvY29sIGluIEtXaW4gR05PTUUgNDggez4=")}

\begin{remark}
  {\ATHENA} docking uses \tmverbatim{xdg-toplevel-drag-v1}. KDE Plasma 6.1 and
  later support the merged protocol in KWin; GNOME 48 and later support it in
  Mutter. Older Wayland compositors may not support detachable-window dragging
  and redocking correctly. The reference system uses Plasma 6.7.4.
\end{remark}

\INLINE_COMMENT{ATHENA-DATA cmd="object" val=("156", "PGxhYmVsfHJlbWFyazpkb2NraW5nIHVzZXMgeGRndG9wbGV2ZWxkcmFndjEgS0RFIFBsYXNtYSA2MSBhbmQgbGF0ZXIgc3VwcG9ydAp0aGUgbWVyZ2VkIHByb3RvY29sIGluIEtXaW4gR05PTUUgNDggfT4=")}

\INLINE_COMMENT{ATHENA-DATA cmd="object" val=("72", "PGxhYmVsfEgyIFNldHRpbmcgdXAgQVRIRU5BIFRyYW5zbWl0dGVyIG9uIGEgc2VydmVyPg==")}

\subsection{Setting up {\ATHENA} Transmitter on a server}

A typical deployment runs a headless {\ATHENA} backend on loopback port 8765
and \tmverbatim{athena-transmitter} in front of it on port 8766.

Generate the backend delegation keypair and start the backend:
\begin{tmcode}
ATHENA/bin/ATHENA.bin -H \textbackslash
  --delegation-key-dir /var/lib/athena-backend/delegation \textbackslash
  --generate-server-keypair
ATHENA/bin/ATHENA.bin -H \textbackslash
  --rag-server /srv/athena/vault \textbackslash
  --rag-listen-address 127.0.0.1 \textbackslash
  --delegation-key-dir /var/lib/athena-backend/delegation
\end{tmcode}
Create \tmverbatim{/etc/athena-transmitter/config.json}:
\begin{tmcode}
\{
  "listen_address": "127.0.0.1",
  "listen_port": 8766,
  "key_dir": "/var/lib/athena-transmitter",
  "accepted_clients": "/var/lib/athena-transmitter/accepted-clients.json",
  "pending_clients": "/var/lib/athena-transmitter/pending-clients.json",
  "upstream": \{
    "url": "http://127.0.0.1:8765",
    "public_key": "BACKEND_PUBLIC_KEY",
    "fingerprint": "BACKEND_FINGERPRINT"
  \}
\}
\end{tmcode}
Generate the transmitter keypair:
\begin{tmcode}
ATHENA/bin/athena-transmitter --config /etc/athena-transmitter/config.json --generate-keypair
\end{tmcode}
Add trusted client public keys to the transmitter
\tmverbatim{accepted-clients.json}:
\begin{tmcode}
\{
  "accepted": [
    "CLIENT_PUBLIC_KEY"
  ]
\}
\end{tmcode}
Add the transmitter public key to the backend
\tmverbatim{accepted-clients.json} with role \tmverbatim{proxy}:
\begin{tmcode}
\{
  "accepted": [
    \{
      "public_key": "TRANSMITTER_PUBLIC_KEY",
      "role": "proxy"
    \}
  ]
\}
\end{tmcode}
Start the transmitter:
\begin{tmcode}
ATHENA/bin/athena-transmitter --config /etc/athena-transmitter/config.json
\end{tmcode}
Reference systemd units, the optional XClarity Controller power hook, and a
Cloudflare Tunnel unit are provided under
\tmverbatim{tools/athena-transmitter/deploy}.

\INLINE_COMMENT{ATHENA-DATA cmd="object" val=("76", "PGxhYmVsfEgyIFNldHRpbmcgdXAgV2ViLUFjY2Vzc2libGUgQVRIRU5BIG9uIGEgc2VydmVyPg==")}

\subsection{Setting up Web-Accessible {\ATHENA} on a server}

Build the release AppImage as in 2.5, then build the Web-Accessible {\ATHENA}
OCI image:
\begin{tmcode}
tools/athena-web-server/build-image.sh \textbackslash
  --runtime container_build/ATHENA-rel.AppImage \textbackslash
  --image localhost/athena-web:latest
\end{tmcode}
Run the broker as an unprivileged user with rootless Podman and delegated
cgroup v2 memory and PID controllers. On a systemd host, enable memory
accounting for the broker user slice before starting it:
\begin{tmcode}
sudo systemctl set-property user-$(id -u).slice MemoryAccounting=yes
\end{tmcode}
A basic loopback broker invocation is:
\begin{tmcode}
ATHENA/bin/athena-web-server \textbackslash
  --listen-address 127.0.0.1 \textbackslash
  --port 8090 \textbackslash
  --image localhost/athena-web:latest \textbackslash
  --max-connections 4 \textbackslash
  --max-memory 4GiB \textbackslash
  --storage-limit 2GiB
\end{tmcode}
Put a WebSocket-capable TLS reverse proxy or tunnel in front of the broker.
Configure STUN/TURN for the WebRTC media path as required by the deployment.
The reference production templates under
\tmverbatim{tools/athena-web-server/deploy} use a dedicated openSUSE MicroOS
KVM guest for the broker and rootless session containers, with Cloudflare
Tunnel for signaling and Cloudflare Realtime TURN for media relay.

\INLINE_COMMENT{ATHENA-DATA cmd="object" val=("28", "PGxhYmVsfEgyIE5leHQgc3RlcHM+")}

\subsection{Next steps}

Start {\ATHENA} and read the documentation available from the Help menu. The
documentation is still under construction. Additional information is available
at \href{https://athena.evalisk.org/}{athena.evalisk.org}. For questions or
build problems, contact
\href{mailto:nuaptan@evalisk.org}{nuaptan@evalisk.org}.

\end{document}
